\chapter{Instalación paso a paso}\label{cap:instalacion}

\section{Obtener el código fuente}
\begin{enumerate}
    \item Descomprima el paquete del proyecto en su carpeta de documentos o bien abra una terminal (PowerShell) para clonarlo con Git desde el repositorio oficial:
    \begin{lstlisting}[style=consola]
cd C:\Users\SuUsuario\Documents
git clone https://github.com/whosstixi19/AppCuria.git
\end{lstlisting}
    \item Verifique que dentro de la carpeta del proyecto se encuentre la subcarpeta \texttt{AplicacionCuria} (con el archivo \texttt{AplicacionCuria.csproj}) y en la raíz el script de base de datos \texttt{BD.sql}.
\end{enumerate}

\section{Configurar pgAdmin 4 y crear las tablas}

\subsection{Crear la base de datos \texttt{curia}}
\begin{enumerate}
    \item Abra \textbf{pgAdmin 4} en la máquina principal y expanda en el panel izquierdo \textbf{Servers} $\rightarrow$ \textbf{PostgreSQL 18}.
    \item Haga clic derecho sobre \textbf{Databases} y seleccione \textbf{Create} $\rightarrow$ \textbf{Database...}.
    \item En el campo \textbf{Database}, escriba exactamente en minúsculas:
    \begin{lstlisting}[style=consola]
curia
\end{lstlisting}
    \item Haga clic en el botón \textbf{Save} para crear la base de datos.
\end{enumerate}

\subsection{Ejecutar el script de base de datos (\texttt{BD.sql})}
\begin{enumerate}
    \item En el panel izquierdo de pgAdmin 4, haga clic derecho sobre la base de datos \texttt{curia} recién creada y seleccione \textbf{Query Tool} (Herramienta de consulta).
    \item En la barra superior del editor de consultas, haga clic en el ícono de \textbf{Open File} (Abrir archivo) y seleccione el archivo \texttt{BD.sql} (o copie y pegue su contenido completo en el editor).
    \item Asegúrese de no tener sombreada solo una parte del texto y presione el botón \textbf{Execute script} (ícono de reproducción $\blacktriangleright$ o tecla \textbf{F5}).
    \item En la pestaña inferior \textbf{Messages} deberá mostrarse el mensaje \texttt{Query returned successfully}. Con esto quedarán creadas las 24 tablas del sistema y la cuenta de usuario administrador inicial.
\end{enumerate}

\begin{figure}[H]
    \centering
    \includegraphics[width=0.88\textwidth]{imagenes/01-pgadmin-bd.png}
    \caption{Creación de la base de datos \texttt{curia} y ejecución del script \texttt{BD.sql} en pgAdmin 4.}
    \label{fig:pgadmin-bd}
\end{figure}

\begin{alerta}
El repositorio contiene archivos históricos de versiones previas (\texttt{Tablas.sql} y \texttt{TablasNuevo.sql}). Para la instalación utilice únicamente el archivo actualizado \texttt{BD.sql}, ya que está 100\% sincronizado con todos los modelos de datos actuales del sistema. Asimismo, verifique en pgAdmin 4 que no quede seleccionado solo un fragmento de texto antes de presionar \textbf{F5}, para que se ejecute el script completo.
\end{alerta}

\section{Configurar la cadena de conexión}
Por motivos de seguridad, al clonar el proyecto desde GitHub el archivo \texttt{appsettings.json} viene con la cadena de conexión en blanco (\texttt{"DefaultConnection": ""}) y el archivo \texttt{appsettings.Development.json} no se incluye en el repositorio público (está listado en \texttt{.gitignore}).

\begin{alerta}
Si intenta ejecutar \texttt{dotnet run} sin antes configurar la cadena de conexión, la consola mostrará el error \texttt{System.InvalidOperationException: The ConnectionString property has not been initialized}. Además, recuerde que si existe el archivo \texttt{appsettings.Development.json}, ASP.NET Core le da prioridad sobre \texttt{appsettings.json}.
\end{alerta}

\begin{enumerate}
    \item Abra el archivo \texttt{appsettings.json} (y, si existe en su carpeta, también \texttt{appsettings.Development.json}) dentro de \texttt{AplicacionCuria}.
    \item En la sección \texttt{ConnectionStrings}, coloque la cadena de conexión apuntando a la base de datos \texttt{curia} con la contraseña asignada a PostgreSQL (\texttt{Clave\_curia}):
\begin{lstlisting}[style=json]
{
  "ConnectionStrings": {
    "DefaultConnection": "Host=localhost;Port=5432;Database=curia;Username=postgres;Password=Clave_curia"
  }
}
\end{lstlisting}
    \item Guarde los cambios antes de iniciar el servidor.
\end{enumerate}

\section{Ejecutar la aplicación y habilitar el acceso en red local}

\subsection{Ejecución local en la máquina principal}
\begin{enumerate}
    \item Abra una terminal de PowerShell dentro de la carpeta \texttt{AplicacionCuria} (donde se ubica \texttt{AplicacionCuria.csproj}).
    \item Restaure los paquetes del proyecto y ejecute la aplicación:
    \begin{lstlisting}[style=consola]
dotnet restore
dotnet run
\end{lstlisting}
    \item La consola mostrará que el servidor está escuchando en \url{http://localhost:5155} y que la verificación inicial de recordatorios se ejecutó con \texttt{[OK]}.
\end{enumerate}

\subsection{Habilitar el acceso desde otras computadoras de la red (Cliente--Servidor)}
Para que las demás computadoras de la Curia puedan conectarse mediante su navegador web a la máquina principal:

\begin{enumerate}
    \item \textbf{Política de cookies en red local (\texttt{Program.cs}):} Verifique que en la línea 41 del archivo \texttt{Program.cs} la política de cookies esté definida como:
    \begin{lstlisting}[style=codigo]
options.Cookie.HttpOnly = true;
options.Cookie.SecurePolicy = CookieSecurePolicy.SameAsRequest;
\end{lstlisting}
    Esto permite que los navegadores de las computadoras clientes mantengan la sesión iniciada al conectarse por dirección IP local mediante \texttt{http://}.
    \item \textbf{Abrir el puerto en el Firewall de Windows:} En la máquina principal, abra una ventana de PowerShell como \textbf{Administrador} y ejecute por única vez el siguiente comando para permitir conexiones entrantes al puerto \texttt{5155}:
    \begin{lstlisting}[style=consola]
New-NetFirewallRule -DisplayName "AplicacionCuria 5155" -Direction Inbound -LocalPort 5155 -Protocol TCP -Action Allow
\end{lstlisting}
    \item \textbf{Iniciar el servidor escuchando en toda la red local:} En la carpeta \texttt{AplicacionCuria}, ejecute:
    \begin{lstlisting}[style=consola]
dotnet run --urls "http://0.0.0.0:5155"
\end{lstlisting}
    \item \textbf{Consultar la dirección IP o nombre de equipo de la máquina principal:} En otra ventana de PowerShell de la máquina principal, ejecute:
    \begin{lstlisting}[style=consola]
ipconfig
hostname
\end{lstlisting}
    Anote la dirección \textbf{IPv4} (por ejemplo, \texttt{192.168.1.100} o \texttt{172.29.x.x}) o el nombre del equipo (\texttt{hostname}).
    \item \textbf{Ingresar desde las demás computadoras:} En cualquiera de las computadoras clientes de la red, abra el navegador web e ingrese a \texttt{http://IP\_DE\_LA\_MAQUINA\_PRINCIPAL:5155}.
\end{enumerate}

\begin{nota}
El comodín \texttt{0.0.0.0} se utiliza exclusivamente en el comando \texttt{dotnet run} de la máquina principal para indicarle que acepte conexiones por cualquiera de sus tarjetas de red. En el navegador de las demás computadoras siempre debe escribirse la dirección IPv4 real o el nombre de red de la máquina principal.
\end{nota}

\begin{figure}[H]
    \centering
    \includegraphics[width=0.88\textwidth]{imagenes/02-consola-dotnet.png}
    \caption{Ejecución del servidor en PowerShell escuchando en \texttt{http://0.0.0.0:5155}.}
    \label{fig:consola-dotnet}
\end{figure}

\section{Verificación final y credenciales iniciales de acceso}
Al acceder a la dirección desde el navegador web (tanto en la máquina principal como desde un equipo cliente), se mostrará la pantalla de \textbf{Inicio de Sesión} de \textbf{ECLESOFT 2.0}. El script \texttt{BD.sql} incluye precreada la cuenta de administrador lista para el primer ingreso:

\begin{itemize}
    \item \textbf{Correo electrónico:} \texttt{admin@curia.com}
    \item \textbf{Contraseña:} \texttt{Clave\_curia}
\end{itemize}

\begin{figure}[H]
    \centering
    \includegraphics[width=0.85\textwidth]{imagenes/03-login.png}
    \caption{Pantalla de Inicio de Sesión de ECLESOFT 2.0.}
    \label{fig:login}
\end{figure}
